\section{Independent-walk encoding and local quantum surfaces}
\label{tree:sec:background}

We now construct the local passage laws for spanning-tree maps, using
independent walk increments and the disk topology of $\SLE_8$ cells.
These model-specific inputs will feed into
Theorem~\ref{common:identification}. We first record the exact
discrete sewing and the continuum facts needed to construct its local
passage laws. Throughout the spanning-tree argument,
\begin{equation}\label{tree:eq:parameters}
 \gamma=\sqrt2,\qquad \kappa'=\frac{16}{\gamma^2}=8,\qquad
 Q=\frac2\gamma+\frac\gamma2=\frac3{\sqrt2}>2,
 \qquad p=\frac1{d_\gamma},\quad \xi=\gamma p,
 \quad s=\frac p2.
\end{equation}
The number $d_\gamma>2$ is the volume-growth exponent in
Equation~\eqref{tree:eq:background_volume} and the dimension parameter
of the LQG metric. Fix an integer $k$ with $ks>8$.

\subsection{The graph encoded by an axial walk}

Let $Z^{\mathrm{disc}}=(L^{\mathrm{disc}},R^{\mathrm{disc}})$ be a
two-sided random walk on $\Z^2$, with independent increments uniform
on $(1,0),(-1,0),(0,1),(0,-1)$ and $Z^{\mathrm{disc}}_0=0$.
Write $L_i,R_i$ for its coordinates in this subsection. For $i\leq j$, set
\begin{equation}\label{tree:eq:discrete_tree_equivalence}
 i\sim_L j\quad\Longleftrightarrow\quad
 L_i=L_j=\min_{i\leq r\leq j}L_r.
\end{equation}
This is an equivalence relation. Its classes are the primal tree
vertices; a horizontal step joins its two consecutive classes by
a tree edge. The up- and down-traversals of one tree edge are
identified. A vertical up-step with terminal time $u$ is matched
to the vertical down-step with terminal time $v>u$ when
\begin{equation}\label{tree:eq:discrete_R_match}
 R_{u-1}=R_v<\min_{u\leq r<v}R_r.
\end{equation}
Each match adds a unit edge between the primal vertices represented
by $u-1$ and $v-1$. Since a vertical step has constant $L$, either
endpoint of that step represents the same primal vertex. The added
edges may be loops or parallel edges.

\begin{theorem}[Mullin encoding and its infinite version]
\label{tree:thm:background_mullin}
Rooted spherical map--spanning-tree pairs with $n$ edges are in
bijection with length-$2n$ axial walks starting at the origin,
remaining in $\Z_{\geq0}^2$, and ending at the origin. The
reconstruction is \eqref{tree:eq:discrete_tree_equivalence}--
\eqref{tree:eq:discrete_R_match}, with the contour order supplying the
embedding and root. Thus uniform pairs correspond exactly to
the walk conditioned on this excursion event. Applying the
bilateral sewing gives the uniform infinite spanning-tree
decorated planar map $\mathcal M$.
\end{theorem}

The finite assertion, including the root and all non-tree edges,
is Mullin's bijection in the explicit form
\cite[Section~2.3, Theorem~1]{DHBijection}; see also
\cite[Section~3.2, Theorem~3.3]{HSSurvey}. The bilateral
construction and the coordinate-contour identification are
\cite[Sections~3.1.1--3.1.3, Proposition~3.3]{GwynneHoldenSunGraphDistances}.
Each excursion has probability $4^{-2n}$ before conditioning,
which proves the asserted uniformity. A pair has
$1+\#\{\text{first-coordinate up-steps}\}$ vertices: its
spanning tree has exactly one edge for each such up-step.

For a finite integer interval $I=[a,b]$, define its
\emph{internal graph} $G_I$ by the same rules, retaining every
horizontal step in $I$, including unmatched traversals, and
only vertical matches whose two steps belong to $I$.
Let $\pi_I$ be the projection of contour times to its vertices,
and let $d_I$ be its all-edge graph distance. Both are functions
of $(Z^{\mathrm{disc}}_t-Z^{\mathrm{disc}}_a)_{a\leq t\leq b}$.
The equivalence relation inside $I$ is the restriction of the
bilateral relation: its minimum test already concerns exactly
the interval between the two times. Hence $G_I$ maps injectively
on vertices to $\mathcal M$, with every edge a genuine map edge.
It is generally not an induced subgraph. If $I\subset J$, every
path in $G_I$ is available in $G_J$. Consecutive contour vertices
are connected at cost at most one, so $G_I$ is connected.

The only discrete metric theorem imported here is the following
consequence of \cite[Theorem~1.6, model~2]{DingGwynneDimension}: for the
root vertex $\rho$ of $\mathcal M$,
\begin{equation}\label{tree:eq:background_volume}
 \lim_{R\to\infty}\frac{\log\#B_{d_{\mathcal M}}(\rho,R)}
 {\log R}=d_\gamma\quad\hbox{almost surely}.
\end{equation}
In particular, for each $\varepsilon>0$, the probability that
this ball has more than $R^{d_\gamma+\varepsilon}$ vertices tends
to zero. These distances use all primal edges.

\subsection{Quantum surfaces, Brownian contours, and interval cells}

A quantum surface is an equivalence class of fields in conformal
charts, where a conformal map $\phi:U\to V$ changes a field $h$
on $V$ to $h\circ\phi+Q\log|\phi'|$ on $U$. Its area measure
is denoted by $\mu_h$ and its quantum boundary length by $\nu_h$.
The constant-shift conventions are
\begin{equation}\label{tree:eq:measure_shift}
 \mu_{h+c}=e^{\gamma c}\mu_h,\qquad
 \nu_{h+c}=e^{\gamma c/2}\nu_h.
\end{equation}
The ordinary unit-area sphere is the sphere measure of weight
$4-\gamma^2$ disintegrated at $\mu_h(\mathcal S)=1$, as in
\cite[Definition~4.21(ii)]{DMS}. Its distinguished points are
forgotten at the end. By \cite[Proposition~A.13]{DMS}, these
points are conditional quantum-area samples, so forgetting them
gives exactly the ordinary unmarked sphere law.

\begin{theorem}[Peanosphere and interval surfaces]
\label{tree:thm:background_peanosphere}
The following statements hold at $\gamma=\sqrt2$.
\begin{enumerate}
\item A $\gamma$-quantum cone decorated by an independent
space-filling $\SLE_8$ curve $\eta$, subsequently parameterized
by $\mu_h$-area with $\eta(0)=0$, has a bilateral boundary-length
process $Z=(L,R)$ with independent Brownian coordinates. The
process determines the decorated quantum surface up to its
conformal coordinate equivalence.
\item For the ordinary unit-area sphere, the corresponding
area-parameterized curve has duration one, with
$\eta(0)=\eta(1)$, and its boundary-length process consists of
two independent Brownian excursions. It determines the
decorated sphere.
\item On the cone, for a deterministic finite interval $[a,b]$, the restricted
increments $(Z_t-Z_a)_{a\leq t\leq b}$ determine the abstract
quantum surface $\eta([a,b])$, its restricted curve, and its
incoming and outgoing boundary arcs with their quantum lengths.
The determination is modulo conformal maps. Abstract decorated
surfaces from disjoint bilateral intervals are independent.
For the sphere, the same reconstruction holds on
$[\delta,1-\delta]$ for every fixed $0<\delta<1/2$, with its
excursion increment law.
\end{enumerate}
\end{theorem}

The cone assertions are \cite[Theorems~1.9 and~1.11]{DMS}; the
interval assertion is also explicitly proved in the proof of
\cite[Corollary~9.3]{DMS}, with independence in the proof of
\cite[Lemma~8.2]{DMS}. A precise formulation of interval
determination is \cite[Lemma~4.9]{GHSSurvey}. The sphere assertion
is \cite[Theorems~1.1 and~3.1]{MillerSheffieldSphere}. The symmetric-interval
reconstruction is proved in the final paragraph of its proof of
Theorem~1.1 for $\gamma\leq\sqrt2$, on p.~49. The correlation is
$-\cos(\pi\gamma^2/4)=0$. A fixed deterministic normalization of
contour heights matches the Brownian variance convention to the
walk variance $1/2$ per coordinate per step; this choice affects
neither the quantum area parameter nor the target metric law.
For completeness, bilateral contour convergence follows by
the ordinary central limit theorem for bounded independent
increments: their covariance is $\tfrac12 I_2$, and disjoint
time increments remain independent. The bound
$\E|Z^{\mathrm{disc}}_b-Z^{\mathrm{disc}}_a|^4\leq C|b-a|^2$
gives tightness of the linearly interpolated paths on bounded
windows. The excursion version and
its density estimates are addressed in Section~\ref{tree:sec:sphere}.

Here are the precise topological facts used when passing from
contours to cells. Two Brownian times have the same image if and
only if they are identified by the equivalence relation
generated by
\begin{equation}\label{tree:eq:Brownian_identification}
 L_t=L_u=\inf_{[t,u]}L
 \quad\hbox{or}\quad
 R_t=R_u=\inf_{[t,u]}R,\qquad t\leq u.
\end{equation}
Away from the common sphere endpoint, every class is a singleton,
a pair, or a triple, and every nonsingleton class belongs to just
one of the two coordinates. In a triple $t_1<t_2<t_3$, $t_2$
is a strict local minimum of that coordinate; different strict
local minima have different heights. These statements are the
Brownian identification description in
\cite[Section~1.3.1, Figure~1.6 and footnote~2,
Proposition~1.6, Lemma~8.15]{DMS}.
The mixed-class exclusion includes both temporal orientations, as
stated in the paragraph accompanying Figure~1.6 of~\cite{DMS}.
Its footnote explains the record-set argument: locally these sets
have the law of zero sets of reflected Brownian motions, with time
reversal for the opposite orientations, and a point is polar for
the corresponding planar Brownian motion. Brownian minima on rational
intervals are unique, and distinct strict local minima have distinct
heights. Four scalar-equivalent times would give two such minima at
the same height. Local absolute continuity transfers these facts
between finite Brownian windows and excursion interiors.

These fiber facts supply the contact identifications. We will also
need paths through a cell's interior. For each fixed $a<b$, the
cell $H_{a,b}:=\eta([a,b])$ is
homeomorphic to a closed disk when $\kappa'=8$
\cite[Section~4.1.2]{DingGwynneDimension}. Thus its interior is connected
and both tips are accessible from the interior. We use this
simultaneously for a prescribed countable collection of
intervals and for independent Poisson endpoints. Every
deterministic $t\in(a,b)$ has $\eta(t)\in H_{a,b}^{\circ}$
almost surely, since lying on its frontier would impose a
one-sided record-minimum condition at that deterministic time.
Moreover,
\begin{equation}\label{tree:eq:cell_interior_preimages}
 \eta^{-1}(H_{a,b}^{\circ})\subset(a,b).
\end{equation}
Indeed an outside visit to an interior point would, by
continuity, put the image of a positive-length outside interval
inside $H_{a,b}^{\circ}$. This contradicts area parametrization,
which makes images of time intervals with disjoint interiors
overlap only in zero quantum area. The same argument places
both tips on the boundary. The curve is proper in the bilateral
plane: $\eta(t)\to\infty$ as $t\to\pm\infty$. Hence the
preimage of a planar compact set is bounded. None of these
topological assertions identifies an internal metric limit at
a cell boundary.

\subsection{The established LQG metric}

\begin{theorem}[Reference metric]
\label{tree:thm:background_metric}
With one deterministic normalization fixed, the $\gamma$-LQG
metric $D_h$ is a measurable function of $h$, is a length metric,
and induces the conformal-chart topology. Its internal metric
$D_h(\cdot,\cdot;U)$ is determined by $h|_U$. For continuous $f$,
\begin{equation}\label{tree:eq:metric_weyl}
 D_{h+f}(x,y)=\inf_{P:x\to y}
   \int_0^{\len(P;D_h)}e^{\xi f(P(t))}\,dt,
\end{equation}
where the infimum is over $D_h$-rectifiable paths parameterized
by $D_h$-length. For conformal \mbox{$\phi:U\to V$},
\begin{equation}\label{tree:eq:metric_covariance}
 D_{h\circ\phi+Q\log|\phi'|}(x,y;U)
 =D_h(\phi(x),\phi(y);V).
\end{equation}
On the whole plane geodesics exist between all points, and for
each fixed pair the geodesic is almost surely unique. The metric
extends intrinsically to the cone and sphere; the unit-area
sphere is almost surely a nondegenerate compact metric space.
\end{theorem}

Existence and the length, locality, and Weyl properties are
\cite[Theorems~1.1--1.2 and Section~1.2]{GwynneMillerMetric}; geodesic and
local-extension conventions are recorded in its Sections~1.2 and~1.5,
including Remark~1.5. Full conformal covariance is
\cite[Theorem~1.3]{GwynneMillerConformal}. In particular, for
$c_v=\gamma^{-1}\log v$,
\begin{equation}\label{tree:eq:reference_quantum_unit}
 D_{h+c_v}=v^pD_h.
\end{equation}
Only the already established metric is assigned these axioms.

We also use internal control near artificial flat faces. If
$K$ is compactly contained in a field chart and
$0<\alpha<\xi(Q-2)$, then almost surely there are finite random
$C$ and $r_0>0$ such that every sufficiently small dyadic open
square $S$ meeting $K$, with closure in the chart and side
length $r<r_0$, satisfies
\begin{equation}\label{tree:eq:internal_square_holder}
 \sup_{x,y\in S}D_h(x,y;S)\leq Cr^\alpha.
\end{equation}
This follows from the simultaneous internal-square bound
\cite[Lemma~3.20, Equation~(3.67)]{DFGPSWeakMetric}, using dyadic
Borel--Cantelli; see also its Proposition~3.10. It controls
approach to a flat face from within the square. Local absolute
continuity transfers this almost-sure assertion to the charts
used here. Quantitative probability bounds for a family of
changes of field require separate uniform estimates.

\subsection{Local fields and locally retained Poisson cells}

Use $h$ for the real quantum field and $\widehat h$ for the
independent imaginary-geometry field generating $\eta$ before
area parametrization. Its angular period is $2\pi\chi$, where
$\chi=2/\sqrt\kappa-\sqrt\kappa/2=1/\sqrt2$ and
$\kappa=16/\kappa'=2$; see
\cite[Section~2.1.3]{GwynneMillerSheffieldHarmonic} for this convention.
Precisely, take a whole-plane GFF $\widehat h_0$ whose average
on $\partial\mathbb D$ is zero, independently of $h$, and a
uniform $\Theta\in[0,2\pi\chi)$ independently of both fields.
Then $\widehat h$ is the class of $\widehat h_0+\Theta$ modulo
additive constants in $2\pi\chi\Z$. This is the standard
whole-plane GFF modulo a period
\cite[Section~2.2.1, p.~48]{MillerSheffieldIGIV}.

This phase convention makes the imaginary-field law invariant
under every fixed positive dilation. Indeed, for $r>0$ the
field $\widehat h_0(r\,\cdot)$, after subtracting its circle
average on $\partial\mathbb D$, is a circle-average-zero
whole-plane GFF. Conditional on $\widehat h_0$, its removed
average added to $\Theta$, modulo $2\pi\chi$, is still uniform.
Thus $\widehat h(r\,\cdot)$ has the same law as $\widehat h$.
The domain Markov property in this convention is
\cite[Proposition~2.8]{MillerSheffieldIGIV}.

We specify the joint cone normalization used for changes of
quantum units. For deterministic $v>0$, let
$c_v=\gamma^{-1}\log v$ and use the positive dilation
\[
 r_v=\sup\{r>0:h_r(0)+Q\log r=c_v\},
\]
where $h_r(0)$ denotes the circle average of the cone field.
The circle-average construction makes $r_v\in(0,\infty)$
almost surely and measurable from $h$ alone. Define
\begin{equation}\label{tree:eq:background_joint_zoom}
 \begin{aligned}
  h^{(v)}&=h(r_v\,\cdot)+Q\log r_v-c_v,
  &\widehat h^{(v)}&=\widehat h(r_v\,\cdot),\\
  \eta^{(v)}(t)&=r_v^{-1}\eta(vt),
  &Z^{(v)}(t)&=v^{-1/2}Z(vt).
 \end{aligned}
\end{equation}
By \cite[Proposition~4.13(i)]{DMS}, $h^{(v)}$ has the original
circle-average cone law. The same normalization and curve
transport are written explicitly in
\cite[proof of Theorem~4.4]{GwynneMillerSheffieldHarmonic}.
Conditional on $h$, dilation invariance and independence give
the original imaginary-field law to $\widehat h^{(v)}$.
Consequently $\widehat h^{(v)}$ is independent of $h^{(v)}$.
The dilation is positive, so the imaginary-geometry coordinate
change has no argument term. Its curve is therefore exactly
$\eta^{(v)}$ after area parametrization. Boundary-length scaling
gives $Z^{(v)}$, and the metric transforms as
\[
 D_{h^{(v)}}(x,y)=v^{-p}D_h(r_vx,r_vy).
\]
It follows that the full joint marginal
$(Z^{(v)},h^{(v)},\widehat h^{(v)},\eta^{(v)},D_{h^{(v)}})$ is
the same for every fixed $v$. Countably many fixed values of $v$
use these transformations of one realization. This convention
concerns the continuum background; it does not impose a
scaling rule on an additional passage law.

\begin{proposition}[Background locality]
\label{tree:prop:background_locality}
The following properties concern the fields, curve, and marks.
\begin{enumerate}
\item On disjoint buffered subdomains of a GFF chart, conditioning
on the exterior gives independent zero-boundary GFF parts and
exterior-measurable harmonic parts, for each of the two fields.
For changes of the real reference field, fix one pinning circle
outside all the subdomains and their dependency buffers; the
imaginary constant is interpreted modulo its angular period.
On compact patches inside the cone's unit disk and away from zero,
the actual normalized real-field law is locally equivalent to this
GFF restriction. The sphere has the corresponding local absolute
continuity modulo constants for background properties invariant under
constant shifts. Its metric law will instead be transferred using
the excursion encoding in Section~\ref{tree:sec:sphere}.
\item For an open $U$ and a positive spatial buffer contained
in a chart, $\widehat h$ on that buffer determines all oriented
curve segments between entrance into and exit from $U$.
The real field gives their quantum-area parametrization and
the quantum lengths of their directed frontier arcs.
\item Add a Poisson point process of area times of intensity
$\varepsilon^{-1}$, independent of the surface and curve.
The cells between consecutive marks which are compactly
contained in $U$, with their directions and frontier lengths,
are determined by the two local fields and the spatial marks
in $U$. Conditional on the two fields, the spatial marks form
a Poisson process of intensity $\varepsilon^{-1}\mu_h$ and
have independent restrictions to disjoint open sets. Nested
Poisson refinements may be used with their level labels.
\end{enumerate}
\end{proposition}

The GFF decomposition is \cite[Lemmas~2.1--2.2]{GwynneMillerSheffieldHarmonic};
the buffer in the curve assertion is exactly that of
\cite[Lemma~2.4]{GwynneMillerSheffieldHarmonic}. The retained-cell statement is
\cite[Lemma~3.7]{ContrerasHipGwynneCurvature}, with the spatial marked-point
data specified there. The following elementary details explain
the form used here. By the mapping theorem for Poisson
processes and area parametrization, the conditional intensity
of the spatial marks is $\varepsilon^{-1}\mu_h$, regardless of
the curve. Thus their restrictions are conditionally
independent. Locally oriented curve segments locate consecutive
marks whenever the entire intervening cell stays in $U$;
compact containment permits the stated buffer and exhaustion.
Restriction of the area and frontier-length measures then
supplies the local parametrizations and directed lengths.
Independent additional Poisson processes give nested
refinements; retaining their level labels preserves locality.
These facts specify local cell data, not a conditional law for
graph distances.

For the asserted transfer of field laws, in the circle-average
embedding the cone field on the unit disk has the law of a
whole-plane GFF minus $\gamma\log|\cdot|$, with the same circle
average normalization
\cite[Definition~4.10]{DMS}; see
\cite[Section~4.1.1]{DingGwynneDimension}. On a buffered patch compactly
contained in $\mathbb D\setminus\{0\}$, this logarithm is smooth.
Extend it with a smooth cutoff supported away from the pinning circle.
The Cameron--Martin theorem and the GFF domain comparison
\cite[Lemma~2.1]{DingGwynneDimension} then give local equivalence for
the actual field restrictions. Reference tests use such patches or
their fixed affine images, with one fixed pinning circle outside the
full dependency region. Intrinsic coordinate and quantum-unit identities
remain available on other domains, but do not assert absolute continuity
under random reanchoring. For the unit-area sphere,
\cite[Lemma~2.4]{BorgaGwynneSunPermutons} gives local absolute continuity
modulo constants, with smooth logarithmic terms away from the marks.
We use this assertion for constant-invariant background properties,
not to identify an absolutely normalized passage law.

\Needspace{12\baselineskip}
\subsection{The form of confluence used below}

\begin{proposition}[Dense-pair coverage of geodesic subarcs]
\label{tree:prop:background_confluence}
Let $h$ be a whole-plane GFF and fix a countable dense set
$\mathcal Q\subset\C$. Almost surely, every nonempty open
subarc of every nonconstant $D_h$-geodesic contains a
nonconstant subsegment of the geodesic between some two points
of $\mathcal Q$. The same local conclusion holds for geodesic
subarcs in buffered field charts on the cone or sphere, with
internal metrics used for the localization.
\end{proposition}

\begin{proof}
The strong confluence theorem is
\cite[Theorem~2]{BhatiaKavvadiasConfluence}. More specifically, its
Propositions~5--6 say that a strictly interior geodesic segment
is uniquely minimizing, and that sufficiently small
perturbations of its endpoints force every connecting geodesic
to contain a prescribed shorter interior segment. Given an
open subarc of $P:[0,T]\to\C$, choose $s<t$ strictly inside its
parameter interval and $0<\varepsilon<(t-s)/2$. Choose points
of $\mathcal Q$ in the endpoint neighborhoods from
Proposition~6. Their geodesic contains
$P([s+\varepsilon,t-\varepsilon])$. There are countably many
pairs in $\mathcal Q$, so fixed-pair uniqueness holds for all
of them on one event of probability one.

For a local chart first shorten the subarc so that it lies
in a smaller domain with positive metric distance to the
boundary of a larger chart. The internal version of strong
confluence, proved in
\cite[Section~6.1, proof of Proposition~50,
Equation~(93)]{BhatiaKavvadiasConfluence}, applies there. The field
absolute continuity in Proposition~\ref{tree:prop:background_locality}
transfers this local event. By shortening again, connecting
geodesics between nearby endpoints are too short to reach
the chart boundary, so their relevant segments are also
ambient geodesics. Additive constants rescale every length
equally and do not affect the conclusion.
\end{proof}
